\makereporttitle{2023}{noiembrie}

\makereportpreamble

\textbf{Rezumat:} În urma discuțiilor cu dl.~prof.~Băran și pe baza materialelor primite, am început să lucrez la reproducerea rezultatelor din lucrarea „Fingerprints of classical chaotic dynamics in quantum behavior”, publicată în \textit{Romanian Reports in Physics}, volumul 70 din 2018, articolul nr.~105.

În acest scop, am început să mă familiarizez cu limbajul de programare Julia, parcurgând documentația oficială și cursurile introductive din \textit{JuliaAcademy}.

Pentru a-mi aprofunda cunoștințele de calcul numeric, am început să citesc cartea „A First Course in the Numerical Analysis of Differential Equations” a lui Arieh Iserles, ediția a II-a, publicată în 2012 de \textit{Cambridge University Press}.

\vspace{0.5em}

\textbf{Activitate de documentare:}

Pe lângă înțelegerea și interpretarea rezultatelor prezentate în lucrarea ,,Fingerprints of classical chaotic dynamics in quantum behavior'' (citată mai sus), am început să parcurg literatura la care face trimitere, pentru a înțelege mai bine sistemul cuantic studiat, precum și fenomenul fizic și datele experimentale asociate.
    
Am început să parcurg articolele:
\begin{itemize}
    \item „Collective Nuclear Motion and the Unified Model” al lui A.~Bohr și B.~R.~Mottelson, publicat în 1955 de \textit{North-Holland Publishing};
    \item „Semiclassical treatment of quadrupole boson hamiltonians: new results for small and large amplitude regimes” al lui A.~A.~Raduta, V.~Băran și D.S.~Delion, publicat în 1995 în \textit{Nuclear Physics A};
    \item „Regular and irregular features of classical motion described by a quadrupole boson Hamiltonian” a lui V.~Băran, A.~A.~Raduta și D.~S.~Delion, publicat în 1996 în \textit{Physical Review E};
    \item „Classical and quantal chaos described by a fourth order quadrupole boson Hamiltonian” a lui V.~Băran și A.~A.~Raduta, publicat în 1997 în \textit{International Journal of Modern Physics E}.
\end{itemize}

\vspace{0.5em}

\textbf{Cursuri:}
Am participat la mai multe cursuri de la alte direcții de studiu din cadrul Școlii Doctorale:
\begin{itemize}
    \item ,,Interacția radiațiilor ionizante și neionizante cu celulele și țesuturile biologice'', ținut de prof.~dr.~Aurel Popescu;
    \item ,,Propagarea luminii prin sisteme optice liniare'', ținut de prof.~dr.~Daniela Dragoman;
    \item ,,Conceptual approaches to scientific research'' și ,,Sistemul climatic ca sistem hidrodinamic'', ținute de prof.~dr.~Mihai Dima;
    \item ,,Microscopie electronică'', ținut de prof.~dr.~Victor Ciupină;
    \item ,,Magnetismul structurilor nanodimensionale'', ținut de prof.~dr.~Victor Kuncser;
\end{itemize}
De asemenea, am participat la prelegerea ,,De la mecanica cuantică la calculul cuantic și la calculatorul cuantic'', ținută de dl.~prof.~Adrian Tanasa (Université Bordeaux) în cadrul Seminarului de Informatică Teoretică de la FMI-UB. Am parcurs module 6, 7 și 8 din cursul online ,,Quantum Mechanics'', oferit de Georgetown University prin intermediul platformei \textit{edX}, ținut de prof.~dr.~Jim Freericks.

\vspace{0.5em}

\textbf{Observațiile conducătorului științific}:

\vspace{2em}

\makereportpostamble{29 noiembrie 2023}